\section{Testing \& Continuous Delivery}

% --- SLIDE 1: WHY ML TESTING IS DIFFERENT ---
\begin{frame}{Why You Cannot Just Write Unit Tests}

    \centering
    \begin{columns}[T]
        \begin{column}{0.48\textwidth}
            \begin{tcolorbox}[colback=blue!5, colframe=blue!50, title=\small Ordinary software]
                \small
                \begin{itemize}
                    \item The correct output is \textbf{specified in advance}.
                    \item \texttt{assert add(2,2) == 4}
                    \item A test failing means the code is wrong.
                    \item Behaviour changes only when code changes.
                \end{itemize}
            \end{tcolorbox}
        \end{column}
        \begin{column}{0.48\textwidth}
            \begin{tcolorbox}[colback=red!5, colframe=red!50, title=\small Machine learning]
                \small
                \begin{itemize}
                    \item The correct output is \textbf{learned}, and only known statistically.
                    \item \texttt{assert model.predict(x) == ?}
                    \item A metric dropping might mean the code is wrong, the data changed, or the seed
                          differed.
                    \item Behaviour changes when \emph{data} changes, with no commit at all.
                \end{itemize}
            \end{tcolorbox}
        \end{column}
    \end{columns}

    \vspace{0.9em}
    \begin{block}{So you test three separate things}
        \textbf{(1)} the \emph{code} around the model, with ordinary unit tests;
        \textbf{(2)} the \emph{data} going in, with schema and quality checks;
        \textbf{(3)} the \emph{model's behaviour}, with properties and thresholds rather than exact values.
    \end{block}

\end{frame}

% --- SLIDE 2: TESTING DATA ---
\begin{frame}[fragile]{Testing the Data}

    \begin{columns}[T]
        \begin{column}{0.5\textwidth}
            \small
            \textbf{Two layers}
            \begin{itemize}
                \item \textbf{Schema:} columns present, types correct, categories in the allowed set,
                      nullability as declared. This is a hard failure --- stop the pipeline.
                \item \textbf{Quality / statistics:} distributions within expected bounds, null rate below a
                      ceiling, no duplicate keys, freshness within an SLA. This is usually a warning, and a
                      gate on retraining.
            \end{itemize}
            \vspace{0.4em}
            \textbf{Write the schema once} and reuse it at three points: training input, retraining input,
            and serving input. That single artifact is your best defence against training/serving skew.
        \end{column}
        \begin{column}{0.46\textwidth}
            \begin{lstlisting}[language=Python, basicstyle=\ttfamily\tiny]
def validate(df: pd.DataFrame) -> None:
    assert set(SCHEMA) <= set(df.columns), "missing column"
    assert df["age"].between(18, 120).all()
    assert df["contract"].isin(CONTRACTS).all()
    assert df["monthly_charges"].isna().mean() < 0.01
    assert df["id"].is_unique
    assert (NOW - df["event_ts"].max()).days < 2   # freshness
            \end{lstlisting}
            \vspace{0.3em}
            \begin{tcolorbox}[colback=gray!8, colframe=gray!60, boxrule=0.5pt]
                \scriptsize Tools that formalise this: Great Expectations, Pandera, TensorFlow Data
                Validation, \texttt{dbt} tests. Twenty lines of \texttt{assert} statements already capture
                most of the value.
            \end{tcolorbox}
        \end{column}
    \end{columns}

\end{frame}

% --- SLIDE 3: TESTING MODELS ---
\begin{frame}{Testing the Model}

    \centering
    \renewcommand{\arraystretch}{1.3}
    {\scriptsize
    \begin{tabular}{p{3.0cm}|p{8.4cm}}
        \toprule
        \textbf{Test type} & \textbf{What it asserts} \\
        \midrule
        \textbf{Smoke / shape}   & The model loads, accepts one row, returns a probability in $[0,1]$. Catches
                                   90\% of deployment breakage. \\
        \textbf{Minimum functionality} & Cases so easy that failure is unambiguous: a customer who cancelled
                                   last month must score high risk. \\
        \textbf{Invariance}      & Perturbations that must \emph{not} change the prediction: reordering
                                   columns, changing a customer ID, whitespace in a category. \\
        \textbf{Directional}     & Perturbations with a known sign: raising the price should not \emph{lower}
                                   predicted churn. \\
        \textbf{Golden set}      & A fixed, curated set of examples with expected outputs, reviewed by a human
                                   and version-controlled. \\
        \textbf{Regression / slice} & Overall metric above a floor \emph{and} no important segment worse than
                                   the incumbent by more than $\delta$. \\
        \bottomrule
    \end{tabular}}

    \vspace{0.21em}
    \begin{itemize}
        \scriptsize
        \item \textbf{Slice tests matter more than the headline metric.} An overall $+1\%$ that hides
              $-8\%$ for one region is a regression, not an improvement.
        \item {\scriptsize The invariance/directional/minimum-functionality vocabulary is from Ribeiro et al.,
              ``Beyond Accuracy: Behavioral Testing of NLP Models with CheckList,'' ACL 2020,
              \href{https://arxiv.org/abs/2005.04118v1}{arXiv:2005.04118v1} --- the idea transfers directly to
              tabular models.}
    \end{itemize}

\end{frame}

% --- SLIDE 4: ML TEST SCORE ---
\begin{frame}{A Rubric to Audit Yourself: The ML Test Score}

    \begin{block}{Breck et al., IEEE Big Data 2017}
        \textbf{28 actionable tests} for production readiness, grouped into four sections:
        \textbf{Features and Data}, \textbf{Model Development}, \textbf{ML Infrastructure}, and
        \textbf{Monitoring} (7 monitoring tests).
    \end{block}

    \vspace{0.15em}
    \begin{columns}[T]
        \begin{column}{0.5\textwidth}
            \begin{tcolorbox}[colback=green!5, colframe=green!50!black, title=\scriptsize How it is scored]
                \scriptsize
                \begin{itemize}
                    \item \textbf{Half a point} per test executed manually, with the result documented and
                          distributed.
                    \item \textbf{A full point} if a system runs that test automatically and repeatedly.
                    \item Sum within each of the four sections.
                    \item \textbf{The final score is the \emph{minimum} of the four section scores.}
                \end{itemize}
            \end{tcolorbox}
        \end{column}
        \begin{column}{0.46\textwidth}
            \scriptsize
            \textbf{Why the minimum?}
            \begin{itemize}
                \item Because all four areas are necessary. Beautiful model development with zero monitoring
                      scores zero.
                \item It is deliberately hard to game by doing more of what you already enjoy.
            \end{itemize}
            \vspace{0.12em}
            \begin{tcolorbox}[colback=yellow!8, colframe=yellow!50!black, boxrule=0.5pt]
                \scriptsize Run this on your own project this week. Most teams score well under 1 --- and the
                exercise tells you exactly which section to invest in next.
            \end{tcolorbox}
        \end{column}
    \end{columns}

    \vspace{0.12em}
    {\scriptsize Source: Breck, Cai, Nielsen, Salib \& Sculley, ``The ML Test Score: A Rubric for ML Production
    Readiness and Technical Debt Reduction,'' \textit{Proc.\ IEEE Big Data}, 2017.
    \href{https://research.google/pubs/the-ml-test-score-a-rubric-for-ml-production-readiness-and-technical-debt-reduction/}{Paper}.}

\end{frame}

% --- SLIDE 5: CI/CD ---
\begin{frame}{CI/CD for ML: What Runs When}

    \centering
    \renewcommand{\arraystretch}{1.3}
    {\small
    \begin{tabular}{p{3.0cm}|p{4.2cm}|p{4.2cm}}
        \toprule
        \textbf{Trigger} & \textbf{What runs} & \textbf{Gate} \\
        \midrule
        Every commit / PR & Linting, unit tests, a training run on a \emph{tiny} sample &
            Must be green to merge; must finish in minutes \\
        Merge to main     & Build the container, publish it by digest &
            Image builds and the smoke test passes \\
        New data or a schedule & Full training pipeline, with data validation first &
            Candidate beats the incumbent on the frozen evaluation set \\
        Manual promotion  & Register the model version, then staged rollout &
            A human approves --- or an automated policy does, with a rollback plan \\
        \bottomrule
    \end{tabular}}

    \vspace{0.9em}
    \begin{itemize}
        \item \textbf{Keep the per-commit job fast.} A CI job that trains the real model for six hours is a
              CI job people learn to skip.
        \item \textbf{Continuous training} (retraining on fresh data) and \textbf{continuous delivery}
              (shipping new pipeline code) are different loops running at different cadences. Do not
              entangle them.
    \end{itemize}

\end{frame}

% --- SLIDE 6: STAGED ROLLOUT ---
\begin{frame}{Never Flip the Switch: Shadow, Canary, A/B}

    \centering
    \begin{tikzpicture}[
        b/.style={rectangle, draw, rounded corners=2pt, minimum height=0.7cm, text width=1.7cm,
                  align=center, font=\tiny, inner sep=2pt},
        arr/.style={-{Latex[length=1.6mm]}, thick, gray!70!black}
    ]
        % Shadow
        \node[font=\scriptsize\bfseries, anchor=south] at (1.8,1.5) {1. Shadow};
        \node[b, fill=gray!14]  (s0) at (0,0.6)  {traffic};
        \node[b, fill=blue!12]  (s1) at (2.4,1.1) {old model\\\textbf{serves}};
        \node[b, fill=orange!12](s2) at (2.4,0.1) {new model\\\emph{logs only}};
        \draw[arr] (s0) -- (s1); \draw[arr, dashed] (s0) -- (s2);

        % Canary
        \node[font=\scriptsize\bfseries, anchor=south] at (6.0,1.5) {2. Canary};
        \node[b, fill=gray!14]  (c0) at (4.2,0.6)  {traffic};
        \node[b, fill=blue!12]  (c1) at (6.6,1.1)  {old model\\99\%};
        \node[b, fill=orange!12](c2) at (6.6,0.1)  {new model\\1\%};
        \draw[arr] (c0) -- (c1); \draw[arr] (c0) -- (c2);

        % A/B
        \node[font=\scriptsize\bfseries, anchor=south] at (10.2,1.5) {3. A/B test};
        \node[b, fill=gray!14]  (a0) at (8.4,0.6)  {traffic};
        \node[b, fill=blue!12]  (a1) at (10.8,1.1) {A: 50\%};
        \node[b, fill=orange!12](a2) at (10.8,0.1) {B: 50\%};
        \draw[arr] (a0) -- (a1); \draw[arr] (a0) -- (a2);
    \end{tikzpicture}

    \vspace{0.7em}
    \begin{itemize}
        \small
        \item \textbf{Shadow:} the new model sees real traffic and its predictions are logged and compared,
              but never returned. Catches crashes, latency blow-ups and wild disagreement at \textbf{zero user
              risk}. Do this first, always.
        \item \textbf{Canary:} a small slice of real traffic gets the new model. Watch error rate, latency
              and the business metric; ramp $1\% \rightarrow 10\% \rightarrow 50\%$, with an automatic
              rollback rule written down \emph{before} you start.
        \item \textbf{A/B test:} a designed experiment to measure whether the new model actually improves the
              \emph{business} outcome --- which is not the same question as whether its F1 is higher.
        \item \textbf{Rollback must be one command and it must be rehearsed.} The previous image digest and
              the previous model version are both still in the registry: that is what they are for.
    \end{itemize}

\end{frame}
